\documentclass{article}
\usepackage[UTF8]{ctex}
\usepackage{amsmath,amssymb,amsfonts}
\usepackage{enumitem}
\usepackage{geometry}
\geometry{a4paper, margin=1in}
% === 错题本排版优化（不参与题目标识） ===
\setlength{\parindent}{0pt}
\setlength{\parskip}{0.35em plus 0.10em minus 0.08em}
\setlength{\emergencystretch}{3em}
\setlist{itemsep=0.25em, topsep=0.40em, parsep=0pt}
\allowdisplaybreaks[2]
% === 错题本排版优化结束 ===
\begin{document}
\section*{太原理工大学2024年高等代数真题}
\section*{59．太原理工大学2024年高等代数试卷}

\subsection{题目1}
一．（10 分）计算 \(n\) 阶行列式

\[
\left|\begin{array}{cccccc}
x & y & 0 & \cdots & 0 & 0 \\
0 & x & y & \cdots & 0 & 0 \\
\vdots & \vdots & \vdots & & \vdots & \vdots \\
0 & 0 & 0 & \cdots & x & y \\
y & 0 & 0 & \cdots & 0 & x
\end{array}\right|
\]

\subsection{题目2}
二．（20 分）已知线性方程组

\[
\left\{\begin{array}{l}
3 x_{1}+x_{2}-x_{3}+2 x_{4}=2 ; \\
x_{1}-5 x_{2}+2 x_{3}+x_{4}=-1 ; \\
2 x_{1}+6 x_{2}-3 x_{3}-3 x_{4}=a+1 ; \\
-x_{1}-11 x_{2}+5 x_{3}+4 x_{4}=4 .
\end{array}\right. \text { (可能有误) }
\]

（1）a 为何值时线性方程组有解？\\
（2）当有解时，求出它的一般解。

\subsection{题目3}
三。（20 分）解答如下问题：

\begin{itemize}
  \item[（1）] 证明：如果 \((f(x), g(x))=1,(f(x), h(x))=1\) ，那么 \((f(x), g(x) h(x))=1\) 。
  \item[（2）] 证明：如果 \((f(x), g(x))=1\) ，那么 \((f(x) g(x), f(x)+g(x))=1\) ．
\end{itemize}

\subsection{题目4}
四．（20分）方阵 \(A\) 如果满足 \(A^{2}=E\) ，那么称 \(A\) 是对合矩阵。设 \(A, B\) 都是数域 \(P\) 上的 \(n\) 阶方阵，证明：

\begin{itemize}
  \item[（1）] 如果 \(A, B\) 都是对合矩阵，且 \(|A|+|B|=0\) ，那么 \(A+B, E+A B\) 都不可逆．
  \item[（2）] 如果 \(B\) 是对合矩阵，且 \(|B|=-1\) ，那么 \(E+B\) 不可逆．
\end{itemize}

\subsection{题目5}
五．（20 分）证明如下问题：

\begin{itemize}
  \item[（1）] \(n\) 阶实对称矩阵 \(A\) 为正定矩阵的充分必要条件是有 \(n\) 阶可逆实矩阵 \(C\) 使得 \(A=C^{\prime} C\) 。
  \item[（2）] \(n\) 阶实对称矩阵 \(A\) 为正定矩阵的充分必要条件是有 \(n\) 阶可逆实对称矩阵 \(C\) 使得 \(A=C^{2}\) 。
\end{itemize}

\subsection{题目6}
原题拆分片段：
六（20 分）求由向量组 \(\alpha_{1}, \alpha_{2}\) 生成的子空间与由向量组 \(\beta_{1}, \beta_{2}\) 生成的子空间的交的基和维数，其中

\[
\alpha_{1}=(1,2,1,0), \alpha_{2}=(-1,1,1,1) ; \beta_{1}=(2,-1,0,1), \beta_{2}=(1,-1,3,7) .
\]

\subsection{题目7}
原题边界“七．”：（20 分）设 \(\mathbb{R}^{4}\) 是所有 4 维实向量构成的线性空间，定义 \(\mathbb{R}^{4}\) 上的线性变换如下
\[
\varphi: \mathbb{R}^{4} \rightarrow \mathbb{R}^{4},\left(x_{1}, x_{2}, x_{3}, x_{4}\right) \mapsto\left(x_{1}+x_{2}, x_{2}-x_{3}, x_{3}-2 x_{4}, x_{1}+2 x_{4}\right) .
\]

\begin{itemize}
\item[（1）] 求 \(\varphi\) 在基 \(\varepsilon_{1}=(1,0,0,0), \varepsilon_{2}=(0,1,0,0), \varepsilon_{3}=(0,0,1,0), \varepsilon_{4}=(0,0,0,1)\) 下的矩阵。
  \item[（2）] \(\varphi\) 是否为同构映射？说明理由．
\end{itemize}

\subsection{题目6}
八．（20 分）已知二次型 \(f\left(x_{1}, x_{2}, x_{3}, x_{4}\right)=2 x_{1} x_{2}+2 x_{1} x_{3}-2 x_{1} x_{4}-2 x_{2} x_{3}+2 x_{2} x_{4}+2 x_{3} x_{4}\) ．

\begin{itemize}
  \item[（1）] 求一个正交线性替换 \(X=Q Y\) 将 \(f\left(x_{1}, x_{2}, x_{3}, x_{4}\right)\) 化为标准形．
  \item[（2）] 求一个可逆线性替换 \(X=P Y\) 将 \(f\left(x_{1}, x_{2}, x_{3}, x_{4}\right)\) 化为规范形．
\end{itemize}
\end{document}
